% Source ownership:
% - Mainly Arthur/Manel/Oliwier
% - pending final benchmark figures / commentary

\section{OpenMP Parallelization}
\label{sec:openmp}

\subsection{\texttt{compute\_total\_energy}}

The dominant cost in \texttt{compute\_total\_energy} is the
\(O(N^2)\) double loop over non-bonded Lennard-Jones pairs. Each
iteration is independent --- it reads a pair of coordinates, evaluates
\texttt{ljpairenergy}, and adds to a running sum --- which makes it a
natural fit for OpenMP worksharing.

The outer loop is distributed across threads with
\texttt{!\$omp parallel do}, with the inner index, the squared
distance \(r^2\), and the dihedral cosine \texttt{cosphi} declared
\texttt{private} to prevent cross-thread interference. The global
accumulators \texttt{elj} and \texttt{etors} are aggregated safely
using a \texttt{reduction(+:...)} clause.

The two loops inside the subroutine have different iteration structures
and were scheduled accordingly. The LJ pair loop has a triangular
iteration space because the inner bound depends on the outer index, so
\texttt{schedule(dynamic,10)} is used to prevent faster threads from
sitting idle while slower ones finish larger chunks of work. The
torsional loop over the carbon backbone is uniform in length and uses
the default static scheduling, which has lower overhead when the
workload is balanced.

Both parallel regions are wrapped in an \texttt{if(omp\_totalenergy)}
runtime guard, which is set from the input file. This allows the same
binary to run in purely serial mode during development or lightweight
tests without any code changes.

\subsection{\texttt{delta\_energy}}

The \texttt{delta\_energy} subroutine is the true bottleneck of the MC
loop --- it is called once per accepted or rejected move, millions of
times per simulation. The key algorithmic feature here is the dynamic
partitioning of atoms into \texttt{fixedlist} and \texttt{movedlist} at
the start of each call.

Only cross-interactions between fixed and moved atoms can have changed
during a pivot, so the LJ loop processes exclusively those pairs rather
than the full \(O(N^2)\) set. This drastically reduces the amount of
work required for a local move and makes the subroutine much more
sensitive to efficient parallelization.

The two nested loops over fixed and moved form a perfectly rectangular
iteration space, which makes \texttt{collapse(2)} effective: it merges
both loops into a single flat range of \(n_{\mathrm{fixed}}
n_{\mathrm{moved}}\) iterations and distributes them evenly across
threads.

Since pivot moves near the ends of the chain move very few atoms,
spawning threads for only a small number of pairs would cost more than
it saves. The parallel block is therefore gated by
\texttt{if(omp\_deltaenergy .and. nfixed*nmoved > 1000)}, ensuring that
OpenMP is only activated when the workload genuinely justifies the
threading overhead.

\subsection{Benchmarking Methodology}

To quantify the effect of thread count on both energy routines, a
dedicated benchmark program \texttt{mainbenchtotal.f90} was written. It
performs 10,000 back-to-back calls to
\texttt{compute\_total\_energy} for a given chain size and reports the
wall time via \texttt{MPI\_Wtime}.

Two shell scripts automate the full benchmark sweeps.
\texttt{4.runomptotaltest.sh} compiles and runs the benchmark binary
across thread counts of 1, 2, 3, and 4 for chain sizes of 20, 50, 100,
and 500 carbons, writing the results to a CSV file in the benchmark
results folder.

\texttt{5.runompdeltatest.sh} follows the same structure but targets
\texttt{delta\_energy} by running full 1M-step MC simulations via
\texttt{mainparallelreplicas.x} and extracting wall time from the CPU
output files. Both scripts are written for the cluster queue and load
the appropriate MPI module.

The results are plotted by \texttt{plotbenchmarks.py}, which reads both
CSV files and produces per-chain-size scaling plots as PDF figures.

A note must be made regarding the \texttt{delta\_energy} benchmark:
because it is measured inside full Monte Carlo production runs, the
timings are affected by randomness. Different random trajectories lead
to different accepted moves and different moved/fixed partitions, so the
workload is not bitwise identical between runs. For that reason, the
\texttt{delta\_energy} benchmark should be interpreted as a practical
end-to-end timing measurement rather than as a perfectly controlled
microbenchmark.

% Oliwier to add images / results / analysis

\subsection{SIMD Exploration}

An attempt was made to additionally exploit single-core SIMD
vectorization on the \texttt{ljpairenergy} function using the
\texttt{!\$omp declare simd} directive. The idea was to generate a
vectorized variant of the function callable from an
\texttt{!\$omp simd}-annotated inner loop, allowing each thread to
process multiple atom pairs simultaneously using AVX registers.

In practice, activating this required structural changes to
\texttt{delta\_energy} that conflicted with the existing threading
design. The \texttt{collapse(2)} clause and the inner-loop
\texttt{cycle} on \texttt{isexcluded} are not naturally compatible with
SIMD regions, and removing them in order to support vectorization
degraded the thread load balancing that \texttt{collapse(2)} provides.

Given that OpenMP threading was the primary parallelization goal, the
SIMD experiment was abandoned and the directive was removed from the
code. It remains a direction worth revisiting, particularly if the pair
loops are later refactored to use pre-filtered pair lists that eliminate
the exclusion branch entirely.